% Chapter 2 -- Understanding the Desktop (Part I, Learning; AN-14 task form).
% Sources (Law 1): os/flair/event.c (one button: the cooker reads button 0
% only; Ctrl is the command modifier); os/flair/finder_desktop.h
% (FINDER_DBLCLICK_TICKS 20 = a third of a second, FINDER_DBLCLICK_SLOP 4 px,
% FINDER_DRAG_SLOP 3 px); os/flair/finder_desktop.c (the volume icon can be
% dragged, the Trash cannot); os/flair/menu.c (press-drag-release tracking,
% dragging across titles, MenuInfo_item_selectable: dimmed items and
% dividers cannot be chosen; MenuKey: the first enabled item with the key;
% the "^N" drawing of a command key); os/milton/kmain.c (band 1 is the
% shell's own bar, band 2 the foreground application's; the Apple menu
% rebuilt from the disk each time it is pulled down; flair_live_about);
% os/flair/apple_menu.c, spec/apple_menu_items.h (About This Computer...,
% then each accessory whose program is on the disk); os/flair/shell_menus.h
% (the upper bar's File, Edit, View and Special menus).  Model: Macintosh
% User's Guide 1991 Ch. 2 "Understanding the Desktop" (apple-macintosh.md
% Sec 8.2).
% GATES: test-flair-desktop-icons test-flair-mouse test-flair-menu test-flair-menu-crossdrag test-menu test-apple-menu test-flair-apple-menu test-finder-menu test-flair-finder-cmds test-pubs-flair-screens
\Gates{test-flair-desktop-icons, test-flair-mouse, test-flair-menu,
  test-flair-menu-crossdrag, test-menu, test-apple-menu,
  test-flair-apple-menu, test-finder-menu, test-flair-finder-cmds,
  test-pubs-flair-screens}
\chapter{Understanding the Desktop}
\label{ch:desktop}

In this chapter you learn the four things you do with the mouse ---
pointing, clicking, double-clicking and dragging --- and how to choose
commands from menus.  Work through the tasks in order; each one takes a
minute or two.

\section{Pointing and clicking}
\index{pointing}\index{clicking}
To \term{point} at something, move the mouse until the tip of the pointer
is on it.  To click it, point at it and then press and quickly release the
mouse button.
\begin{TaskSteps}
\Step Point at the \texttt{INITECH} icon and click it.
\Result{The icon and its name are drawn in reverse: the icon is
selected.}
\end{TaskSteps}

\FlairScreen[scale=4, label=fig:selected]{selected}{The INITECH icon,
selected}

\begin{TaskSteps}[resume]
\Step Click a bare part of the desktop.
\Result{The icon is no longer selected.  Clicking the desktop also makes
the Finder the active application, as you will see later in this
chapter.}
\end{TaskSteps}
Selecting something does not change it.  It marks the thing you want the
next command to act on.

\section{Double-clicking}
\index{double-clicking}
\term{Double-clicking} is clicking twice in quick succession without moving
the mouse.  You double-click an icon to open it.
\begin{TaskSteps}
\Step Point at the \texttt{INITECH} icon and double-click it.
\Result{A window opens, showing what is on the data disk.  You learn about
windows in Chapter~\ref{ch:windows}.}
\end{TaskSteps}
The two clicks must come within a third of a second of each other, and the
pointer must move no more than a few dots between them.  If nothing opens,
you clicked too slowly, or the mouse moved: try again a little faster, and
hold the mouse still.  FLAIR treats two slow clicks as two separate
clicks.

\section{Dragging}
\index{dragging}
To \term{drag} something, point at it, hold down the mouse button, move the
mouse, and release the button.  You drag icons to move them, and you drag
windows by their title bars (Chapter~\ref{ch:windows}).
\Operation{FINDER-ICON-DRAG}
\begin{TaskSteps}
\Step Point at the \texttt{INITECH} icon and hold down the mouse button.
\Step Move the mouse.
\Result{An outline of the icon follows the pointer once the pointer has
moved a few dots.  The icon itself stays where it was until you release
the button.}
\Step Release the button somewhere on the desktop.
\Result{The icon moves to where you released it.  FLAIR remembers the new
place and shows the icon there each time it starts.}
\end{TaskSteps}
The Trash always stays in its place at the lower right of the screen; you
cannot drag it.  To put something in the Trash you drag the thing onto the
Trash (Chapter~\ref{ch:trash}).

\section{Choosing a command from a menu}
\index{menus}\index{choosing a command}
Most of the work you do in FLAIR is done with commands, and commands are
listed in menus.  The menu titles are the words in the menu bars at the top
of the screen.
\begin{TaskSteps}
\Step Click a bare part of the desktop.
\Result{The lower menu bar shows the Finder's menu titles:
\menuname{File}, \menuname{Edit}, \menuname{View}, \menuname{Special} and
\menuname{Help}.}
\Step Point at \menuname{File} in the lower menu bar and hold down the
mouse button.
\Result{The File menu drops down and stays down while you hold the
button.}
\end{TaskSteps}

\FlairScreen[scale=4, label=fig:filemenu]{filemenu}{The Finder's File menu,
with nothing selected}

\begin{TaskSteps}[resume]
\Step Still holding down the button, move the pointer down the menu.
\Result{Each command you pass over is highlighted.  Commands drawn in
grey are not available now; they are not highlighted and cannot be
chosen.}
\Step Move the pointer off the menu and release the button.
\Result{The menu closes and nothing happens.  Releasing the button on a
highlighted command chooses that command.}
\end{TaskSteps}
While you hold down the button you can move the pointer along the menu bar:
each menu you pass drops down in turn.

\SideHead{Dimmed commands}A command is drawn in grey, or \term{dimmed},
when it cannot be used at the moment.  In Figure~\ref{fig:filemenu},
\menuitem{Open} is dimmed because no folder is selected, and
\menuitem{Close Window} because no window of the Finder is open.  Some
commands are dimmed at all times in this version of FLAIR; Chapter~\ref{ch:menus} tells
you which.

\SideHead{Keyboard equivalents}A command with a caret and a letter at its
right, such as \CmdKey{N} beside \menuitem{New Folder}, can also be given
from the keyboard: hold down \key{Ctrl} and press the letter.  It works only
when the command is not dimmed, and only for the menus of the active
application.  You can use small or capital letters.

\section{The two menu bars}
\index{menu bars}
FLAIR has two menu bars, one above the other.

\SideHead{The upper bar}The upper menu bar belongs to FLAIR.  It is the
same whatever you are doing.  At its left is the Apple menu, then the
titles \menuname{File}, \menuname{Edit}, \menuname{View} and
\menuname{Special}, and at its right the clock.

\FlairScreen[scale=4, label=fig:band1file]{band1file}{The File menu of the
upper menu bar}

The four menus of the upper bar each hold the same two commands,
\menuitem{About} and \menuitem{Quit}.  The commands of the Finder, which
you use to manage your files, are in the lower bar.

\SideHead{The lower bar}The lower menu bar belongs to the \term{active
application}: the application whose window is in front, or the Finder when
you have clicked the desktop or a Finder window.  When you bring another
application's window to the front, the lower bar changes to show that
application's menus.  When FLAIR starts, the HELLO window is in front, so
the lower bar holds HELLO's menus.

\section{The Apple menu}
\index{Apple menu}
The Apple menu, at the left of the upper menu bar, is always there.  It
gives you information about the computer and opens the accessories
installed on the data disk.
\begin{TaskSteps}
\Step Point at the apple at the left end of the upper menu bar and hold
down the mouse button.
\Result{The Apple menu drops down.}
\end{TaskSteps}

\FlairScreen[scale=4, label=fig:applemenu]{applemenu}{The Apple menu}

\SideHead{About This Computer}The first command describes your computer.
\SideHead{Accessories}Below the line is one command for each accessory
whose program is on the data disk.  On the data disk as it is delivered
there is one, the \menuitem{Calculator}.  If you remove an accessory's
program from the disk, its command disappears from the menu the next time
you pull the menu down.

\begin{TaskSteps}[resume]
\Step Still holding down the button, move the pointer to \menuitem{About
This Computer...} and release the button.
\Result{An alert box shows the version of InitechDOS, the memory built into
the computer, and the largest block of that memory not in use.}
\end{TaskSteps}

\FlairScreen[label=fig:about]{about}{About This Computer}
\Operation{APPLE-ABOUT}

\begin{TaskSteps}[resume]
\Step Click \texttt{OK}, or press \key{Enter}.
\Result{The alert box closes.}
\end{TaskSteps}

\section{What you have learned}
\begin{itemize}
\item To select an icon, click it.  To open it, double-click it.
\item To move an icon, drag it.
\item To choose a command, point at the menu title, hold down the mouse
  button, move to the command and release the button.
\item Dimmed commands cannot be chosen.  A command with a keyboard
  equivalent can also be given with \key{Ctrl} and a letter.
\item The upper menu bar is FLAIR's own; the lower menu bar belongs to the
  active application.
\end{itemize}
